% Replace the Introduction sections ``Internship context'' and ``Problem
% statement'', then replace Chapter 2 with the following text.
% Add references to the official SIAGE and INSIGHT project presentations in the
% bibliography before submitting the final thesis.

\section{Internship context}

Population ageing raises interconnected challenges for housing, mobility, access
to healthcare, social participation, digital inclusion, and support for caregivers.
These challenges are addressed not only through national policy frameworks, but
also through locally designed initiatives carried by municipalities, associations,
social-protection organisations, care providers, and citizen groups. Such
initiatives constitute situated forms of knowledge: they document how actors in a
specific territory identify local needs and develop responses to them.

This internship was carried out within the International Chair ``Inclusive
Societies and Ageing'' (SIAGE) at the University of Lorraine and contributes to
the INSIGHT research project. SIAGE brings together research in the social
sciences, public-policy studies, and technical disciplines in order to examine the
conditions of inclusion and exclusion associated with ageing. Its work is concerned
with both the diversity of older adults' life courses and the role of social,
institutional, physical, and digital environments in supporting participation,
autonomy, and citizenship.

Among its participatory action-research activities, the Chair develops the
\textit{Senior Citizen Design Workshops}. Developed in France since 2018 in
collaboration with local authorities and the international REIACTIS network, this
approach aims to involve older adults as citizens rather than only as users or
beneficiaries of public action. The workshops are organised around four connected
stages: participatory diagnosis, citizen-centred co-design, dialogue with public
decision-makers, and follow-up and evaluation. They produce action sheets and
other textual material in which participants describe local problems, proposed
solutions, actors, and territorial resources.

The documentary material available to the project extends beyond the action sheets
produced within the Senior Citizen Design Workshops. Numerous initiatives are also
carried locally by other territorial and gerontological actors. Their websites,
reports, project descriptions, and practical sheets contain valuable, localised
knowledge about responses to ageing-related challenges. However, this knowledge is
dispersed across organisations, territories, formats, and editorial practices.

Taken together, these resources form a large but heterogeneous corpus. They are
partially structured and difficult to mobilise systematically for researchers,
local practitioners, policy-makers, and older citizens who may wish to identify or
draw inspiration from initiatives developed in other territories. The purpose of
\textit{Index\_Insight} is to make this dispersed material more visible,
comparable, and reusable while preserving its documentary provenance.

\section{Problem statement}

Although the corpus contains rich descriptions of local responses to population
ageing, the absence of a shared structure limits its practical and scientific use.
It is difficult to search across territories, compare initiative types, identify the
actors involved, or examine the themes and needs addressed by the documented
initiatives. This scarcity of accessible and comparable information limits the
capacity of researchers to analyse local ageing policies, of practitioners and local
elected representatives to identify transferable experiences, and of older citizens
to learn from initiatives developed elsewhere.

The problem addressed in this internship is therefore not only the technical
extraction of information from text. It is the design of a documented and
interdisciplinary workflow capable of transforming heterogeneous documentary
resources into structured metadata that can support exploration, comparison,
capitalisation of experience, and the circulation of knowledge between territories.

\chapter{General Context of the Project}

\section{Presentation of the host institution}

The International Chair ``Inclusive Societies and Ageing'' (SIAGE) is a research
initiative of the University of Lorraine dedicated to the study of population ageing
and inclusive societies. It approaches ageing as a multidimensional social issue:
inclusion is shaped by individual life courses, social inequalities, public policies,
professional practices, and the design of physical and digital environments. The
Chair is particularly attentive to the ways in which ageing can affect access to
rights, services, mobility, social participation, and citizenship.

SIAGE develops interdisciplinary research that connects social sciences with other
fields, including public health, engineering, and digital technologies. Its objective
is not to treat technical systems as substitutes for social-science interpretation,
but to use them to organise, compare, and make visible knowledge that can inform
research, public action, and local practice. The Chair also seeks to create links
between researchers, public decision-makers, territorial actors, professionals, and
civil-society organisations.

The Senior Citizen Design Workshops are one of the Chair's participatory
action-research activities. They provide a methodological framework through which
older adults can contribute to the identification of local issues and the
co-construction of responses with territorial actors. The action sheets produced
through these workshops are therefore both documentary resources and traces of a
citizen-centred process of knowledge production.

\section{Presentation of the project}

The internship is situated at the intersection of the INSIGHT project, the SIAGE
Chair, and the Senior Citizen Design Workshops. INSIGHT supports an
interdisciplinary approach to language data by connecting Natural Language
Processing and Artificial Intelligence with research in the Humanities and Social
Sciences. Within this context, \textit{Index\_Insight} is a specific application of
this interdisciplinarity to the field of territorial adaptation to population ageing.

The objective of \textit{Index\_Insight} is not limited to creating a searchable
collection of documents. The project aims to capitalise territorial innovations,
make them visible and comparable, facilitate the transfer of experience between
territories, and support the different stages of the Senior Citizen Design approach:
participatory diagnosis, co-design, dialogue with public decision-makers, and the
capitalisation of experience. In this sense, the Index is intended as a tool for the
production and circulation of knowledge that can be used by researchers, local
authorities, practitioners, and older citizens.

\section{Data used}

The corpus combines two complementary forms of documentary material. First, it
includes action sheets produced by older adults themselves within Senior Citizen
Design Workshops implemented by SIAGE. These documents describe problems,
proposed actions, target publics, local resources, and implementation conditions
identified through participatory work.

Second, the corpus includes descriptions of initiatives carried by institutional,
associative, and territorial actors involved in the adaptation of territories to
population ageing in France. The selection of these sources was conducted in
dialogue with the social-science researchers of the SIAGE Chair. Sources were
retained when their material was relevant to the objectives of the Index and when
their websites or documents contained sufficiently substantial information to be
structured. The detailed composition of the final corpus and the reasons for
retaining or excluding individual sources are presented later in the methodology
chapter.

\section{Constraints and needs}

The heterogeneity of the corpus does not result only from technical differences
between HTML pages, PDF documents, and plain text. It also reflects the
multi-level organisation of the French gerontological field. Public institutions,
local authorities, associations, social-protection organisations, care providers,
and other actors intervene at national, regional, departmental, intermunicipal, and
local levels. They pursue complementary objectives relating to prevention,
participation, housing, mobility, health, digital inclusion, and social connection.

This institutional and territorial diversity produces a naturally dispersed
documentary landscape, with variable document formats, levels of detail,
vocabularies, and ways of presenting initiatives. Relevant information may be
distributed across several pages of the same website or located in different
sections of a report. The resulting technical challenge is therefore inseparable
from the fragmentation of the field itself. It justifies the collection,
harmonisation, and structuring work undertaken in \textit{Index\_Insight}.
